\begin{enumerate}

\item Las incógnitas son $R_A$, $M_A$ y $R_B$, y hay tres ecuaciones: equilibrio vertical, momentos y flector nulo en $H$. La rótula hace isostática una viga que sin ella sería hiperestática.

\item La parte $HC$ solo contiene $R_B$ como incógnita, así que el flector nulo en $H$, escrito con la parte derecha, da $R_B$ directamente.

\item Con $R_B$ conocida, el equilibrio vertical da $R_A$ y el de momentos respecto a $A$ da el momento del empotramiento.

\item Los tramos son $AH$, $HD$, $DB$ y $BC$. En $AH$ y $HD$ la ley es la misma, porque en la rótula no actúa ninguna carga; el cortante salta en $D$ y en $B$.

\item El cortante del tramo $DB$ no se anula, de modo que el flector máximo positivo está en $D$, donde el cortante cambia de signo; el mínimo está en el empotramiento.

\end{enumerate}

\noindent\textbf{Operaciones y resultados intermedios.}

\begin{enumerate}

\item Momento de la carga de HB respecto a H.

\[\fitmath{y_{1}=\frac{\frac{46}{7}\cdot {4}^{2}}{2}+\frac{\frac{6}{7}\cdot {4}^{3}}{3}\approx 70.85714\,\mathrm{kN m}}\]

\item Reacción en B (flector nulo en H).

\[\fitmath{y_{2}=\frac{\frac{496}{7}+15\cdot 2-6\cdot 6}{4}\approx 16.21429\,\mathrm{kN}}\]

\item Reacción en A.

\[\fitmath{y_{3}=\frac{4+10}{2}\cdot 7+15-6-\frac{\frac{496}{7}+15\cdot 2-6\cdot 6}{4}\approx 41.78571\,\mathrm{kN}}\]

\item Flector en el empotramiento, con la parte derecha.

\[\fitmath{y_{4}=\frac{\frac{496}{7}+15\cdot 2-6\cdot 6}{4}\cdot 7+6\cdot 9-\left(\frac{4\cdot {7}^{2}}{2}+\frac{\frac{6}{7}\cdot {7}^{3}}{3}\right)-15\cdot 5\approx -103.5\,\mathrm{kN m}}\]

\end{enumerate}
